
\subsubsection{2.1 KV Cache Compression}

KV cache eviction has been studied extensively as sequence lengths grow. H2O \citep{Zhangetal2023} uses accumulated attention scores as a heuristic for token importance. StreamingLLM \citep{Xiaoetal2024} retains "attention sink" tokens plus a sliding window. Locret \citep{Huangetal2024} trains lightweight two-layer MLP retaining heads per transformer layer on a language modeling objective, learning a contextual importance score (CIS) for each token; it achieves up to 20× compression with limited perplexity increase. PruLong \citep{Princeton2024} learns binary attention head masks. More recent work addresses systems-level concerns including quantization-eviction interactions and memory hierarchy optimization [Gokhale et al., 2025; Yao et al., 2025; Chu et al., 2025].

A shared characteristic of all these eviction methods is that they are trained with language modeling loss. Their eviction priorities are thus calibrated for next-token prediction. When applied to models that have been task-fine-tuned, the eviction policy retains tokens relevant to language modeling, while the fine-tuned model's representations have been reorganized around task-discriminative features. This misalignment is the motivation for JointLoRA-KV.

\subsubsection{2.2 Parameter-Efficient Fine-Tuning}

LoRA \citep{Huetal2022} injects low-rank matrices A∈ℝ^{d×r}, B∈ℝ^{r×d} (with r ≪ d) into selected linear projections, enabling task adaptation with a small fraction of trainable parameters. QLoRA \citep{Dettmersetal2023} extends this to quantized base models. Surveys of PEFT methods are provided in Mao et al. [2024] and Yang et al. [2024]. In existing deployment pipelines, PEFT fine-tuning and KV compression are applied sequentially: the model is adapted first, then a fixed compression policy is layered on top. No prior work had jointly optimized these two components for a task-specific objective prior to this work.

\subsubsection{2.3 Joint PEFT and KV Compression}

The closest prior approach is arXiv:2604.21335 [Jiang \& Wang, 2026], which combines LoRA routing paths with KV value-group routing in a single architecture. This work establishes that LoRA and KV routing components can coexist in the forward pass. However, it trains with language modeling loss throughout — not task classification loss — and evaluates on perplexity and RULER benchmarks rather than GLUE or LongBench. The amazon-science/icr-kv-caching work [Molfese et al., EACL 2026] analyzes post-hoc how fine-tuning affects KV compression robustness but does not propose joint training. The survey by Kim et al. [2025] explicitly identifies the absence of unified compression-plus-tuning frameworks as an open problem; JointLoRA-KV is a direct response to this gap.

\subsubsection{2.4 Differentiable Sparse Selection}

The soft-to-hard training pattern in JointLoRA-KV — sigmoid relaxation during training, hard discrete selection at inference — is established in quantization-aware training and sparse learning. Locret itself uses soft scoring during training. Bengio et al. [2013] introduced the straight-through estimator for propagating gradients through discrete operations. JointLoRA-KV applies this pattern to the joint LoRA + KV eviction training problem with a task classification objective.

